
% This one will format for two-sided binding (ie left and right pages have mirror margins; blank pages inserted where needed):
\documentclass[a4paper,twoside]{gamosinos}


%----------------------------------------------------------------------------------------
%	NOMBRE DEL AUTOR
%----------------------------------------------------------------------------------------

%% HERE CHANGE NAME
\fancyfoot[C]{\emph{Name author, \today}}  %HERE


%----------------------------------------------------------------------------------------
%	PACKAGES
%----------------------------------------------------------------------------------------

\usepackage{pdfpages}

\usepackage{listings}
\usepackage{ctable}
\usepackage{xcolor}
\usepackage{xcolor,colortbl}


\usepackage{algorithm}
\usepackage{algorithmic}

\usepackage{mdframed}

\definecolor{codegreen}{rgb}{0,0.6,0}
\definecolor{codegray}{rgb}{0.5,0.5,0.5}
\definecolor{codepurple}{rgb}{0.58,0,0.82}
\definecolor{backcolour}{rgb}{0.95,0.95,0.92}

\lstdefinestyle{mystyle}{
    backgroundcolor=\color{backcolour},   
    commentstyle=\color{codegreen},
    keywordstyle=\color{magenta},
    numberstyle=\tiny\color{codegray},
    stringstyle=\color{codepurple},
    basicstyle=\ttfamily\footnotesize,
    breakatwhitespace=false,         
    breaklines=true,                 
    captionpos=b,                    
    keepspaces=true,                 
    numbers=left,                    
    numbersep=5pt,                  
    showspaces=false,                
    showstringspaces=false,
    showtabs=false,                  
    tabsize=2
}
\lstset{style=mystyle}

%%
%% Julia definition (c) 2014 Jubobs
%%
\lstdefinelanguage{Julia}%
  {morekeywords={abstract,break,case,catch,const,continue,do,else,elseif,%
      end,export,false,for,function,immutable,import,importall,if,in,%
      macro,module,otherwise,quote,return,switch,true,try,type,typealias,%
      using,while},
   sensitive=true,
   alsoother={},
   morecomment=[l]\#,
   morecomment=[n]{\#=}{=\#},%
   morestring=[s]{"}{"},%
   morestring=[m]{'}{'},%
}[keywords,comments,strings]%

\lstset{%
    language         = Julia,
    basicstyle       = \ttfamily,
    keywordstyle     = \bfseries\color{blue},
    stringstyle      = \color{magenta},
    commentstyle     = \color{codegreen},
    showstringspaces = false,
}

\lstdefinestyle{Matlab}{  
  belowcaptionskip=1\baselineskip,
  breakatwhitespace=false,        % sets if automatic breaks should only happen at whitespace
  breaklines=true,                % sets automatic line breaking
  xleftmargin=\parindent,
  language=C,
  morekeywords={end, elseif},
  comment=[l]\%,                              % comments    
  tabsize=4, 
  numberstyle=\tiny\noncopynumber,
  showstringspaces=false,
  basicstyle=\footnotesize\ttfamily,
  keywordstyle=\color[rgb]{0,0,1},
  commentstyle=\itshape\color[rgb]{.133,.545,.133},
  stringstyle=\color{mauve},
  directivestyle=\bfseries\color{purple},
  framexleftmargin=3mm,
  frame=single,
}

\lstset{%
    language         = Matlab,
    basicstyle       = \ttfamily,
    keywordstyle     = \bfseries\color{blue},
    stringstyle      = \color{magenta},
    commentstyle     = \color{codegreen},
    showstringspaces = false,
}


%%
%% OpenSCAD definition
%%
\lstdefinelanguage{OpenSCAD}%
  {keywords={sphere,cylinder,acos,atan2,translate,rotate,norm},
      morekeywords=[2]{module,false,true},
  keywordstyle=[2]\color{cyan},
   sensitive=true,
   alsoother={},
   morecomment=[l]{//},
   morestring=[s]{"}{"},%
   morestring=[m]{'}{'},%
}[keywords,comments,strings]%

\lstset{%
    language         = OpenSCAD,
    basicstyle       = \ttfamily,
    keywordstyle     = \bfseries\color{blue},
    stringstyle      = \color{magenta},
    commentstyle     = \color{codegreen},
    showstringspaces = false,
}

%%
%% BDF definition
%%
\lstdefinelanguage{BDF}%
  {keywords={ID,SOL,TIME,TITLE,SUBTITLE,LOAD,SPC,DISP,STRESS,STRAIN,ELFORCE,SPCFORCE,BAILOUT,MDLPRM,SPC1,FORCE,GRID,MAT1,CBAR,PBARL,BAR,ALL},
  morekeywords=[2]{CEND,BEGIN,PARAM,BULK,ENDDATA},
  keywordstyle=[2]\color{cyan},
  sensitive=true,
  comment=[l]{\$},
}[keywords,comments,strings]%

\lstset{%
    language         = BDF,
    basicstyle       = \ttfamily,
    keywordstyle     = \bfseries\color{blue},
    stringstyle      = \color{magenta},
    commentstyle     = \color{codegreen},
    showstringspaces = false,
    literate =  {=}{{{\color{red}=}}}1
                {,}{{{\color{red},}}}1,
}

%%
%% VTU definition
%%

\lstdefinelanguage{VTU}{%
  language     = XML,
  morekeywords = {VTKFile,UnstructuredGrid,Piece,Points,DataArray,Cells,CellData},
}

%%
%% STL definition
%%

\lstdefinelanguage{STL}{%
  morekeywords = {solid,endsolid,facet,normal,outer,loop,vertex,endloop,endfacet},
}

%%
%% DOS/Terminal definition
%%

\lstdefinelanguage{DOS}{
    backgroundcolor=\color{black},
    basicstyle=\scriptsize\color{white}\ttfamily,
    numbers=none,
}

\usepackage{afterpage}

\newcommand\blankpage{%
    \null
    \thispagestyle{empty}%
    \addtocounter{page}{-1}%
    \newpage}



\correctionstrue

% To make text superscripts shortcuts
	\renewcommand{\th}{\textsuperscript{th}} % ex: I won 4\th place
	\newcommand{\nd}{\textsuperscript{nd}}
	\renewcommand{\st}{\textsuperscript{st}}
	\newcommand{\rd}{\textsuperscript{rd}}

%----------------------------------------------------------------------------------------
%	START DOCUMENT
%----------------------------------------------------------------------------------------

\begin{document}
\pagenumbering{gobble}

%%%%% TITLE PAGE INFORMATION
% Everybody needs to complete the following:
\title{Suitably impressive title} %HERE
\author{\textbf{Your Name}} % HERE
\college{Nonlinear Solid Mechanics and Structures Group\\
(GAMOSINOS)}
\degree{Aeronautics and Space Engineer}
\degreedate{2022} % HERE


%%%%% CHOOSE YOUR LINE SPACING
% This is the official option.  Use it for your submission copy and library copy:
\setlength{\textbaselineskip}{22pt plus2pt}
% This is closer spacing (about 1.5-spaced) that you might prefer for your personal copies:
%\setlength{\textbaselineskip}{18pt plus2pt minus1pt}

% You can set the spacing here for the roman-numbered pages (acknowledgements, table of contents, etc.)
\setlength{\frontmatterbaselineskip}{17pt plus1pt minus1pt}

% Leave this line alone; it gets things started for the real document.
\setlength{\baselineskip}{\textbaselineskip}

%%%%% CHOOSE YOUR SECTION NUMBERING DEPTH HERE
% You have two choices.  First, how far down are sections numbered?  (Below that, they're named but
% don't get numbers.)  Second, what level of section appears in the table of contents?  These don't have
% to match: you can have numbered sections that don't show up in the ToC, or unnumbered sections that
% do.  Throughout, 0 = chapter; 1 = section; 2 = subsection; 3 = subsubsection, 4 = paragraph...

% The level that gets a number:
\setcounter{secnumdepth}{2}
% The level that shows up in the ToC:
\setcounter{tocdepth}{2}



% JEM: Pages are roman numbered from here, though page numbers are invisible until ToC.  This is in
% keeping with most typesetting conventions.

%\afterpage{\blankpage}
\maketitle

\pagenumbering{arabic}
\begin{romanpages}

%\afterpage{\blankpage}
%----------------------------------------------------------------------------------------
%	QUOTATION PAGE
%----------------------------------------------------------------------------------------

\vspace*{0.1\textheight}
\newpage
\vspace*{0.2\textheight}
\noindent\enquote{\itshape Thanks to my solid academic training, today I can write hundreds of words on virtually any topic without possessing a shred of information, which is how I got a good job in journalism.}\bigbreak

\hfill Dave Barry
\afterpage{\blankpage}
%----------------------------------------------------------------------------------------
%	ACKNOWLEDGEMENTS -- Nothing to do here except comment out if you don't want it.
%----------------------------------------------------------------------------------------


\begin{acknowledgements}
 	\input{text/acknowledgements}
\end{acknowledgements}

%----------------------------------------------------------------------------------------
%	ABSTRACT
%----------------------------------------------------------------------------------------

%%%%%  Nothing to do here except comment out if you don't want it.
\begin{abstract}
	\input{text/abstract}
\end{abstract}

%%%%% MINI TABLES
% This lays the groundwork for per-chapter, mini tables of contents.  Comment the following line
% (and remove \minitoc from the chapter files) if you don't want this.  Un-comment either of the
% next two lines if you want a per-chapter list of figures or tables.
\dominitoc % include a mini table of contents
%\dominitoc  % include a mini list of figures
%\dominilot  % include a mini list of tables
%\dominilol
%\dominilol

% This aligns the bottom of the text of each page.  It generally makes things look better.
\flushbottom

% This is where the whole-document ToC appears:
\tableofcontents

\listoffigures
	\mtcaddchapter
% \mtcaddchapter is needed when adding a non-chapter (but chapter-like) entity to avoid confusing minitoc

% Uncomment to generate a list of tables:
\listoftables
	\mtcaddchapter

\lstlistoflistings
	%\mtcaddchapter
	
\listofalgorithms
	%\mtcaddchapter
	
	%\mtcaddchapter
%%%%% LIST OF ABBREVIATIONS,symbols, constants...
% This example includes a list of abbreviations.  Look at text/abbreviations.tex to see how that file is
% formatted.  The template can handle any kind of list though, so this might be a good place for a
% glossary, etc.
\include{text/symbols}
\include{text/abbreviations}

% The Roman pages, like the Roman Empire, must come to its inevitable close.
\end{romanpages}

%----------------------------------------------------------------------------------------
%	ADD CHAPTERS AND APPENDIX HERE
%----------------------------------------------------------------------------------------

%%%%% CHAPTERS
% Add or remove any chapters you'd like here, by file name (excluding '.tex'):
\flushbottom
\input{text/chapters/list_of_chapters}

% Starts lettered appendices, adds a heading in table of contents, and adds a
%    page that just says "Appendices" to signal the end of your main text.
\startappendices
% Add or remove any appendices you'd like here:
\input{text/appendix/list_of_appendices}


%----------------------------------------------------------------------------------------
%	REFERENCES
%----------------------------------------------------------------------------------------

\bibliographystyle{plain}
\bibliography{references.bib}



\end{document}
